\emph{Setting.} On a spatially flat Robertson-Walker field (Definition \claimref{obj.synthesis.gtg-flat-spacetime-gauge-fields}, the derived metric $a^{2}(\eta)(d\eta^{2} - d\vec x^{2})$) the rescaled field $\chi = a^{3/2}\psi$ obeys the flat Dirac equation with the time-dependent mass $M = ma(\eta)$, and for a plane wave of momentum $\vec k$ and fixed helicity the two components obey $Iu_{1}' = Mu_{1} + ku_{2}$, $Iu_{2}' = ku_{1} - Mu_{2}$, with $I = \sigma_{1}\sigma_{2}\sigma_{3}$ the unit (below) (standard; the mode equations are those of Landete, Navarro-Salas and Torrenti, arXiv 1305.7374, rewritten in conformal time, checked). That is $Iu' = Hu$ with $H = M\sigma_{3} + k\sigma_{1}$ acting on a two-component spinor of the Pauli algebra. The script checks in the Pauli algebra (GA, signature $(+,+,+)$) that $I = \sigma_{1}\sigma_{2}\sigma_{3}$ is central with $I^{2} = -1$, that $H^{2} = k^{2} + M^{2}$, and that $\tfrac12(ab - ba) = I(a \times b)$ on all pairs of basis vectors, hence on all vectors. For $\rho = \tfrac12(1 + s)$ the equation $I\rho' = [H, \rho] = \tfrac12[h, s] = I(h \times s)$ gives $s' = 2h \times s$ with the factor $2$ from $\rho = \tfrac12(1 + s)$ and $[h,s] = 2I(h \times s)$; so the unit $I$, the restriction of $I_{4}$ to the even subalgebra, is the unit of the mode equation. For two pure states, $|\langle\phi|\varphi\rangle|^{2} = \tfrac12(1 + s_{\phi}\cdot s_{\varphi})$, so the Bogoliubov coefficient of a state of Bloch vector $s$ against the pure state $s'$ (the weight of $s$ on the state of Bloch vector $-s'$) is $|\beta|^{2} = \tfrac12(1 - s\cdot s')$.

\emph{(i)} Write $\epsilon = 1/k$, $s = \sum_{j}\epsilon^{j}c_{j}$ with $c_{j}$ vector polynomials in independent jet variables $M_{0}, M_{1}, \dots$, the derivative acting by $M_{j} \mapsto M_{j+1}$. Then $s' = 2h \times s$ with $h = \epsilon^{-1}\sigma_{1} + M\sigma_{3}$ is, at order $\epsilon^{-1}$, $\sigma_{1} \times c_{0} = 0$, so $c_{0} = \pm\sigma_{1}$ (positive frequency: $+$), and at order $\epsilon^{n}$
\[
c_{n}' = 2\sigma_{1} \times c_{n+1} + 2M\sigma_{3} \times c_{n}.
\]
The part of $c_{n+1}$ orthogonal to $\sigma_{1}$ is therefore $-\tfrac12\sigma_{1} \times (c_{n}' - 2M\sigma_{3} \times c_{n})$, provided the bracket is orthogonal to $\sigma_{1}$; its $\sigma_{1}$ part is fixed by $|s|^{2} = 1$ at order $\epsilon^{n+1}$, $2c_{0}\cdot c_{n+1} = -\sum_{a + b = n+1,\ a,b \ge 1}c_{a}\cdot c_{b}$. So the formal solution is unique once $c_{0}$ is chosen. The script builds it to order $9$ and checks at every order the solvability condition, the norm, the equation's residual and the homogeneity of weight; it finds $c_{1} = M\sigma_{3}$ and $c_{2} = -\tfrac12M^{2}\sigma_{1} + \tfrac12M_{1}\sigma_{2}$. Homogeneity holds for all $n$ by induction: $c_{0}$ has weight $0$, the derivative raises the weight by one and multiplication by $M$ (weight $1$) does too, and the norm terms are products of weights summing to $n + 1$. Hence $M_{n}$, of weight $n + 1$, cannot occur in $c_{j}$ for $j \le n$, and in $c_{n+1}$ it can occur only as the monomial $M_{n}$ itself; that monomial cannot come from $M\sigma_{3} \times c_{n}$ (which carries a factor $M$) nor from the norm terms (products of two factors of positive weight), only from $c_{n}'$, where it is the derivative of the monomial $M_{n-1}$ of $c_{n}$. So $\lambda_{n}v_{n} = -\tfrac12\lambda_{n-1}\sigma_{1} \times v_{n-1}$ with $\lambda_{0}v_{0} = \sigma_{3}$; since $\sigma_{1} \times \sigma_{3} = -\sigma_{2}$ and $\sigma_{1} \times \sigma_{2} = \sigma_{3}$, this gives $v_{n}$ alternating between $\sigma_{3}$ and $\sigma_{2}$, $|\lambda_{n}| = 2^{-n}$ and the stated signs; the script confirms the recursion, the axis and the size for $n \le 8$. At $m = 0$ all jet variables vanish and every $c_{j}$, $j \ge 1$, being homogeneous of positive weight, vanishes: the conformal vacuum is exact.

\emph{(ii)} $\hat h = (\sigma_{1} + \epsilon M\sigma_{3})(1 + \epsilon^{2}M^{2})^{-1/2}$ is expanded with exact binomial coefficients. It agrees with $c_{0}$ and $c_{1}$ and has no $\sigma_{2}$ part, so $s_{\mathrm{ad}} - \hat h = \tfrac12M_{1}\sigma_{2}\epsilon^{2} + O(\epsilon^{3})$; the script computes $1 - s_{\mathrm{ad}}\cdot\hat h$ through order $9$ and finds the displayed terms at orders $4$ and $6$, zero at all odd orders and at orders below $4$. Thus $|\beta_{k}|^{2} = M_{1}^{2}/(16k^{4}) + O(k^{-6})$. In the instantaneous vacuum every sector's negative-frequency state is $-\hat h$, with energy $-|h|$; in the adiabatic state it is $-s_{\mathrm{ad}}$, with energy $-|h|\,s_{\mathrm{ad}}\cdot\hat h$. The conformal-frame energy density is $a^{-4}$ times a positive multiple of the mode sum of these energies (the Hamiltonian density of $\chi$, by the field equation), so the adiabatic state exceeds the instantaneous one by a positive multiple of $\int k^{2}|h|(1 - s_{\mathrm{ad}}\cdot\hat h)\,dk$, whose integrand is $M_{1}^{2}/(8k) + O(k^{-3})$: logarithmically divergent where $M_{1} \neq 0$, while the number moment $\int k^{2}|\beta_{k}|^{2}dk$ converges. On each dyadic shell $[K, 2K]$ the leading term of $\int k^{3}|\beta_{k}|^{2}dk$ contributes at least $M_{1}^{2}/32$, a fixed amount, so the divergence is not removed by any finite subtraction. Adiabatic states of infinite order of the Dirac field on Robertson-Walker fields are Hadamard (Hollands; Fact \claimref{synthesis.imported-dirac-hadamard-form-propagates-and-adiabatic-states-are-hadamard}(c)). If the instantaneous vacuum were Hadamard, its two-point functions would differ from those of such a state by a smooth kernel (Fact \claimref{synthesis.imported-microlocal-hadamard-condition}(b)), and the coincidence limit of the energy density of the difference, a differential operator applied to a smooth kernel and evaluated on the diagonal, would be finite; it is the limit of the mode sums cut off smoothly at large $k$, whose integrand is positive at large $k$, so by monotone convergence it is the divergent integral above. Hence, given Hollands' theorem, the instantaneous vacuum is not Hadamard, and its point-split energy density relative to the adiabatic state diverges to $-\infty$ (a non-Hadamard state has no renormalised energy density of its own).

\emph{(ii), deposits diagonal over the instantaneous vacuum.} A gauge-invariant quasi-free state of one sector is fixed by its one-particle density, a positive element $\tfrac12(\tau + b)$ of the sector's Pauli algebra with scalar part $\tau$ and Bloch part $b$; the adiabatic state (negative-frequency state filled) has $\tfrac12(1 - s_{\mathrm{ad}})$. A deposit diagonal over the instantaneous vacuum, with particle and antiparticle occupations $n_{+}$, $n_{-}$, has a density commuting with $h$, so $b = t\hat h$ with $t = n_{+} + n_{-} - 1$. The $\sigma_{2}$ component of the difference of the two densities is $\tfrac12(t\hat h + s_{\mathrm{ad}})\cdot\sigma_{2} = \tfrac12 s_{\mathrm{ad}}\cdot\sigma_{2} = \tfrac14M_{1}k^{-2} + O(k^{-4})$, since $\hat h$ has no $\sigma_{2}$ part at any order; the script checks this for several values of $t$, including $t = 0$ and $t = \pm1$. This gap does not depend on the occupations. If such a deposit were Hadamard, its two-point functions would differ from the adiabatic state's by a smooth kernel; projecting onto one helicity and extracting the $\sigma_{2}$ component (pseudodifferential operators of order zero, pseudolocal, as in (iii) below) leaves a smooth scalar kernel whose symbol has the fixed-sign leading term $\tfrac14M_{1}k^{-2}$, and its coincidence value is then, by monotone convergence, the divergent integral $\int k^{2}\cdot k^{-2}dk$. So no deposit diagonal over the instantaneous vacuum is Hadamard where $M_{1} \neq 0$, given Hollands' theorem.

\emph{(iii)} Let jets agree below order $n$ and differ by $\Delta M_{n}$ at order $n$ (higher derivatives may differ too). By (i), $c_{j}$ for $j \le n$ is unchanged and $c_{n+1}$ changes exactly by $\lambda_{n}\Delta M_{n}v_{n}$; the script substitutes $M_{n} \mapsto M_{n} + D$ and confirms this, and that $1 - s_{L}\cdot s_{R} = D^{2}4^{-n}/2\cdot\epsilon^{2n+2} + O(\epsilon^{2n+3})$ for $2n + 2 \le 9$; for general $n$ it follows from $1 - s_{L}\cdot s_{R} = \tfrac12|s_{L} - s_{R}|^{2}$ and $|s_{L} - s_{R}| = 2^{-n}|\Delta M_{n}|\epsilon^{n+1} + O(\epsilon^{n+2})$. An adiabatic state of infinite order before $\eta_{d}$ has, at $\eta_{d}$, Bloch vectors agreeing with the left series to every order in $1/k$ (that is what infinite order means for these homogeneous states, Hollands loc. cit.), and the state evolves continuously through $\eta_{d}$; the adiabatic states after $\eta_{d}$ agree with the right series. So $|\beta_{k}|^{2} = \tfrac12(1 - s_{L}\cdot s_{R}) = (\Delta M_{n})^{2}/(2k)^{2n+2} + O(k^{-2n-3})$. The moments: $\int k^{2}|\beta_{k}|^{2}dk$ diverges for $n = 0$ and converges for $n \ge 1$; $\int k^{3}|\beta_{k}|^{2}dk$ diverges logarithmically for $n = 1$ and converges for $n \ge 2$. For every finite $n$ the state is not Hadamard: if it were, its two-point functions would differ from those of a right adiabatic state of infinite order by a smooth kernel, whose equal-time Fourier symbol in each helicity sector is $\tfrac12(s_{L} - s_{R})\cdot\sigma$ in the $\vec k$-dependent frame of that sector. Projecting onto one helicity and extracting the $v_{n}$ component are Fourier multipliers homogeneous of degree zero and smooth away from $k = 0$, hence (after a cut-off near $k = 0$, which adds a smooth kernel) pseudodifferential operators of order zero, which are pseudolocal (standard, cited from memory) and keep the kernel smooth. The resulting scalar kernel has a symbol whose leading term $\tfrac12\lambda_{n}\Delta M_{n}k^{-n-1}$ has a fixed sign at large $k$, so applying the Laplacian $j$ times with $2j \ge n - 1$ and evaluating at coincidence gives, by monotone convergence, a divergent integral $\int k^{2 + 2j - n - 1}dk$; a smooth kernel gives a finite value. So, given Hollands' theorem (which supplies the Hadamard comparison state), the state is not Hadamard.

\emph{(iv)} In the setting of Proposition \claimref{synthesis.sudden-descent-feature-not-tilt}, with $z = k\eta$ and units $(2k)^{-1/2} = 1$, the Bunch-Davies mode is $u = e^{-I_{4}z}(1 - I_{4}/z)$ and the instantaneous start has $\chi = 1$, $\chi_{z} = -I_{4}$ at $z_{d}$ (numbers read in $\mathrm{span}\{1, I_{4}\}$). Then $|u|^{2} = 1 + z^{-2}$ and $\tfrac12(|u_{z}|^{2} + |u|^{2}) = 1 + \tfrac12z^{-4}$, against $1$ and $1$ for the instantaneous data; the script checks both identities and $|\beta|^{2} = \tfrac14z_{d}^{-4}$ at eleven rational points, more than the nine that determine Laurent polynomials of degree at most eight in $z^{-1}$. In physical units the per-mode deficits are $(2k)^{-1}(k\eta_{d})^{-2}$ in $\langle\chi^{2}\rangle$ and $\tfrac k2\cdot\tfrac12(k\eta_{d})^{-4}$ in the energy, and with the measure $k^{2}dk/2\pi^{2}$ both mode sums diverge logarithmically. Compare with an infrared-regularised Bunch-Davies state: the Bunch-Davies modes for $k \ge k_{0}$, any regular modes below $k_{0}$ (the massless Bunch-Davies state itself is infrared divergent on slice fields of nonzero mean, Fact \claimref{synthesis.imported-ford-parker-infrared-divergence}(b), (c)). Such a state is Hadamard (standard, cited from memory), the modes below $k_{0}$ adding only a smooth kernel. If the instantaneous start were Hadamard, its two-point function would differ from the comparison state's by a smooth kernel (the smooth-difference property of Hadamard states, stated for the scalar field in Fact \claimref{synthesis.imported-microlocal-hadamard-condition}(a)), so the coincidence value of the difference of $\langle\chi^{2}\rangle$ on the descent slice would be finite. It is the limit of the smoothly cut-off mode sums of the per-mode differences, which for $k \ge k_{0}$ are positive and equal to $(2k)^{-1}(k\eta_{d})^{-2}$, so by monotone convergence it is the divergent integral $\int^{\infty}dk/(2k\eta_{d}^{2})$. So the start is not Hadamard.

\emph{Consequence.} A massive Dirac deposit equal to, or diagonal over, the instantaneous vacuum at an instant where $ma' \neq 0$, or carried through a descent across which $M = ma$ is not smooth, is not Hadamard. The instantaneous vacuum has finitely many quanta above an adiabatic state of infinite order, so the number of quanta does not decide the Hadamard property. Nothing here says that a Hadamard deposit's reference vacuum must be adiabatic: an adiabatic state of infinite order is itself finitely many pair-coherent quanta over the instantaneous vacuum, and the converse of Hollands' theorem is not used. Status of the parts: (i) and the series statements in (ii) to (iv) are proved by exact algebra; the Hadamard conclusions are given Hollands' theorem, the smooth-difference property of Hadamard states, the Hadamard property of the infrared-regularised comparison state in (iv), and the pseudolocality of pseudodifferential operators.
